\section{Análisis comparativo}

La Tabla~\ref{tab:state-of-the-art} permite observar una regularidad importante: la literatura revisada no avanza en una sola dirección, sino en dos ejes parcialmente tensionados. El primero es el de \textbf{efectividad predictiva}. En este eje, los trabajos más fuertes son aquellos que explotan arquitecturas profundas o modelos basados en transformers, como LSTM-CNN, Bi-GRU con modelos preentrenados, TCN con autoatención y modelos ajustados para severidad depresiva \citep{tadesse2020suicide, mirtaheri2024tcn, shuvo2024bigru, qasim2025depression}. Estos enfoques suelen reportar mejores resultados cuantitativos porque modelan dependencias largas, contexto y variación semántica con mayor riqueza que las aproximaciones basadas en n-gramas o léxicos aislados.

El segundo eje es el de \textbf{interpretabilidad clínica}. Desde esta dimensión, destacan menos los clasificadores de alto rendimiento y más los trabajos que reestructuran la tarea para alinearla con marcos clínicos o evidencias observables. \citet{gaur2018dsm5} y \citet{gaur2019severity} incorporan taxonomías médicas y esquemas ordinales útiles para intervención; \citet{bao2025redsm5} introducen anotación a nivel de síntoma y racionales; \citet{singh2024evidence} convierten el problema en extracción y resumen de evidencia; y \citet{bauer2024llm} utilizan LLMs para interpretar dimensiones discursivas, no solo para predecir. En otras palabras, la interpretabilidad mejora cuando el diseño metodológico obliga al sistema a producir salidas más cercanas al lenguaje clínico.

También se observan diferencias en la unidad de análisis. Algunos artículos clasifican publicaciones individuales. Otros modelan usuarios, trayectorias o cambios temporales. Un tercer grupo analiza comunidades completas o subreddits como indicadores diagnósticos indirectos \citep{tadesse2020suicide, chatterjee2022multimodal, mirtaheri2024tcn, dechoudhury2016shifts, gaur2019severity, gaur2018dsm5, bauer2024llm}. Esta diferencia no es menor: los enfoques por publicación tienden a maximizar detección inmediata, mientras que los enfoques por usuario o transición aportan mayor valor para monitoreo temprano y contextualización clínica. Este contraste ayuda a explicar por qué el estado del arte no debe leerse como una secuencia de resúmenes, sino como una comparación de decisiones metodológicas con consecuencias distintas.

Por último, la comparación muestra que la calidad de los datos condiciona fuertemente el tipo de conclusión posible. Los artículos de revisión señalan que gran parte del campo depende de etiquetas heurísticas, diccionarios o autodeclaraciones \citep{montejo2024survey, abdulsalam2024suicidal}. En cambio, los trabajos con mayor potencial clínico invierten más en anotación experta y en estructuras diagnósticas explícitas, pero a costa de datasets más pequeños y costosos de construir \citep{gaur2019severity, bao2025redsm5}. La lección metodológica es clara: no existe una sola frontera del estado del arte; la frontera predictiva y la frontera interpretativa avanzan a ritmos distintos y exigen decisiones de diseño diferentes. Este punto resulta especialmente importante para el reto, porque anticipa que una solución sólida deberá justificar tanto su rendimiento como el tipo de evidencia que ofrece.
